\section*{Bab \ref{ch:b3:hilbert} --- Ruang Hilbert}

\begin{solution}{exo:b3:hilbert:1}
(a) Yang real: uraikan $\norm{x \pm y}^2 = \norm x^2 \pm 2\langle
x,y\rangle + \norm y^2$ lalu kurangkan. Yang kompleks (dengan \omterm{def:b3:hilbert:inner}{hasil kali dalamnya}
linear pada slot keduanya): dengan menguraikannya seperti di atas,
\[
\langle x, y\rangle = \frac14\sum_{k=0}^{3}
\iu^k\,\bigl\|\iu^kx + y\bigr\|^2,
\]
sebab setiap sukunya menyumbang $\iu^k\cdot2\operatorname{Re}\bigl(
(-\iu)^k\langle x,y\rangle\bigr)$, yang jumlahnya $4\langle
x,y\rangle$ (periksalah keempat nilai $k$-nya; sedangkan $\sum\iu^k
(\norm x^2 + \norm y^2) = 0$).

(b) Untuk $L^1$: $f = \mathbf 1_{\intcc0{1/2}}$ dan $g = \mathbf
1_{\intcc{1/2}1}$: maka $\norm{f\pm g}_1^2 = 1$ masing-masing, berjumlah $2$;
padahal $2\norm f_1^2 + 2\norm g_1^2 = 1 \neq 2$. Untuk norma supnya: $f =
\mathbf 1$ dan $g(t) = t$ pada $\intcc01$: maka $\norm{f + g}_\infty^2 +
\norm{f-g}_\infty^2 = 4 + 1 = 5 \neq 4 = 2 + 2$. Karena gagal memenuhi
hukum jajaran genjang, norma ini tak diinduksi \omterm{def:b3:hilbert:inner}{hasil kali dalam} mana pun
(sebab \omterm{def:b3:hilbert:inner}{hasil kali dalam} akan memaksanya lewat penguraian langsung).
\end{solution}

\begin{solution}{exo:b3:hilbert:2}
(a) $p(f) = \bigl(\int_0^1f\bigr)\mathbf 1$: sebab sungguh $f - \int f
\perp$ fungsi konstan (karena $\int(f - \int f)c = 0$). Jadi penghampiran
konstan terbaik atas $f$ dalam kuadrat rata-rata adalah \emph{rata-ratanya} ---
yakni contoh pertama nilai harapan bersyarat
(\cref{ch:b3:probability}).

(b) $p(f) = f\,\mathbf 1_{\intcc{1/2}1}$: sebab selisihnya
$f\mathbf 1_{\intcc0{1/2}}$ ortogonal terhadap setiap $g$ yang
lenyap pada $\intcc0{1/2}$.

(c) $d^2 = \bigl\|x - \tfrac12\bigr\|_2^2 = \int_0^1(x -
\tfrac12)^2\dd x = \tfrac1{12}$: sehingga $d = \frac1{2\sqrt3}$.
\end{solution}

\begin{solution}{exo:b3:hilbert:3}
(a) Kesetaraannya adalah \cref{thm:b3:hilbert:decomposition}
(sebab $\bar F = (F^\perp)^\perp$, dan $\bar F = H \iff F^\perp =
\{0\}$). Contohnya: ruang $F$ berisi barisan berhingga bersifat padat
di $\ell^2$ (lewat pemancungan) dan sejati: sehingga $F^\perp = \{0\}$ padahal $F
\neq \ell^2$ --- jadi untuk subruang yang tak tertutup, $H = F \oplus
F^\perp$ gagal secara terang-terangan (sebab $F \oplus \{0\} \neq H$).

(b) $\abs{\langle x_n, y_n\rangle - \langle x, y\rangle} \leq
\abs{\langle x_n - x, y_n\rangle} + \abs{\langle x, y_n -
y\rangle} \leq \norm{x_n - x}\sup_n\norm{y_n} + \norm
x\,\norm{y_n - y} \to 0$ (sebab barisan konvergen bersifat terbatas).
Ia dipakai: pada \cref{thm:b3:hilbert:parseval}(1) untuk melihat $x - \lim
S_N \perp e_k$, dan pada \cref{thm:b3:hilbert:decomposition}
untuk melihat $F^\perp = \bar F^{\,\perp}$.
\end{solution}

\begin{solution}{exo:b3:hilbert:4}
Berlaku $e_0 = \frac1{\sqrt2}$. Berikutnya, $x \perp \mathbf 1$ sudah sejak awal
(sebab $\int_{-1}^1x = 0$), dan $\int_{-1}^1x^2 = \frac23$: sehingga $e_1 =
\sqrt{\tfrac32}\,x$. Lalu $x^2 - \langle e_0, x^2\rangle e_0 =
x^2 - \frac13$ (dan $\perp e_1$ menurut paritasnya), dengan
\[
\int_{-1}^1\Bigl(x^2 - \frac13\Bigr)^2\dd x = \frac25 -
\frac49 + \frac29 = \frac{8}{45}:
\qquad e_2 = \sqrt{\tfrac{45}8}\,\Bigl(x^2 - \frac13\Bigr).
\]
Perbandingannya: $P_0 = 1$, $P_1 = x$, $P_2 = \frac{3x^2 - 1}2$, dan
$\sqrt{n + \tfrac12}\,P_n$ memberikan $\frac1{\sqrt2}$,
$\sqrt{\frac32}x$, $\sqrt{\frac52}\,\frac{3x^2-1}2 =
\sqrt{\frac{45}8}\bigl(x^2 - \frac13\bigr)$: yakni tepat $e_0, e_1,
e_2$.
\end{solution}

\begin{solution}{exo:b3:hilbert:5}
Untuk $f(t) = t$: $c_0 = 0$ dan, lewat pengintegralan parsial, $c_n =
\frac{\iu(-1)^n}{n}$ untuk $n \neq 0$: sehingga $\abs{c_n}^2 =
\frac1{n^2}$. Lalu \omterm{thm:b3:hilbert:parseval}{Parseval}:
\[
\frac1{2\pi}\int_{-\pi}^\pi t^2\dd t = \frac{\pi^2}3
= \sum_{n\neq0}\frac1{n^2} = 2\sum_{n\geq1}\frac1{n^2}
\ \Longrightarrow\ \sum_{n\geq1}\frac1{n^2} = \frac{\pi^2}6 .
\]
Untuk $f(t) = t^2$: $c_0 = \frac{\pi^2}3$ dan $c_n =
\frac{2(-1)^n}{n^2}$ (bila $n \ne 0$). Lalu \omterm{thm:b3:hilbert:parseval}{Parseval}:
\[
\frac1{2\pi}\int_{-\pi}^{\pi}t^4\dd t = \frac{\pi^4}5
= \frac{\pi^4}9 + \sum_{n\neq0}\frac4{n^4}
\ \Longrightarrow\
\sum_{n\geq1}\frac1{n^4} = \frac18\Bigl(\frac{\pi^4}5 -
\frac{\pi^4}9\Bigr) = \frac{\pi^4}{90} .
\]
Tak diperlukan syarat $\mathcal C^1$ sepenggal-sepenggal apa pun:
sebab \cref{thm:b3:hilbert:fourier} mencakup setiap fungsi $L^2$.
\end{solution}

\begin{solution}{exo:b3:hilbert:6}
(a) $\varphi(f) = \int_0^{1/2}f = \langle\mathbf
1_{\intcc0{1/2}},\ f\rangle$: jadi wakilnya adalah $a =
\mathbf 1_{\intcc0{1/2}}$, dan $\norm\varphi = \norm a_2 =
\frac1{\sqrt2}$ (\cref{thm:b3:hilbert:riesz}).

(b) Ambillah fungsi tenda $f_n$ yang berpuncak $1$ di
$\frac12$ dan berpendukung selebar $\frac2n$: maka $f_n(\tfrac12) = 1$
sedangkan $\norm{f_n}_2^2 \leq \frac2n \to 0$: sehingga tak ada konstanta $C$
yang dapat memberikan $\abs{f(\frac12)} \leq C\norm f_2$. Jadi penilaian titik
tak bermakna di $L^2$ --- sebab unsurnya berupa kelas modulo himpunan
nol --- dan perhitungan ini adalah alasan kuantitatifnya.
\end{solution}

\begin{solution}{exo:b3:hilbert:7}
(a) Misalkan $M = \sup_k\norm{x_k}$. Barisan skalarnya
$(\langle e_n, x_k\rangle)_k$ terbatas oleh $M$: sehingga pemetikan
diagonal melahirkan $x_{k_j}$ dengan $\langle e_n, x_{k_j}\rangle
\to \gamma_n$ untuk setiap $n$. Lalu untuk setiap $N$: $\sum_{n\leq
N}\abs{\gamma_n}^2 = \lim_j\sum_{n\leq N}\abs{\langle e_n,
x_{k_j}\rangle}^2 \leq M^2$ (lewat Bessel), sehingga $(\gamma_n) \in
\ell^2$ dan $x = \sum_n\gamma_ne_n \in H$
(\cref{thm:b3:hilbert:parseval}(3)). Untuk $y \in H$:
\[
\abs{\langle y, x_{k_j} - x\rangle}
\leq \Bigl|\sum_{n\leq N}\overline{c_n(y)}\bigl(\langle e_n,
x_{k_j}\rangle - \gamma_n\bigr)\Bigr|
+ 2M\Bigl(\sum_{n>N}\abs{c_n(y)}^2\Bigr)^{1/2},
\]
dengan memakai uraian $\langle y, z\rangle =
\sum\overline{c_n(y)}c_n(z)$ dan Cauchy--Schwarz pada ekornya;
lalu pilihlah $N$ kemudian $j$: yakni kekonvergenan lemah ke $x$.

(b) Berlaku $\langle y, e_n\rangle = c_n(y) \to 0$ untuk setiap $y$
(lewat ekor $\ell^2$-nya): sehingga $e_n \rightharpoonup 0$, padahal $\norm{e_n} =
1$: jadi normanya tak \omterm{def:b3:topology:continuity}{kontinu} secara lemah. Pada (a): $\norm x^2 =
\sum\abs{\gamma_n}^2 \leq \liminf_j\norm{x_{k_j}}^2$ (lewat potongan
berhingganya dan Bessel lagi): sehingga limit lemahnya hanya dapat kehilangan norma.
\end{solution}

\begin{solution}{exo:b3:hilbert:8}
Untuk $y$ yang tetap, $x \mapsto \langle y, Tx\rangle$ merupakan
fungsional linear \omterm{def:b3:topology:continuity}{kontinu}; sehingga Riesz memberikan tepat satu $T^*y$ dengan
$\langle y, Tx\rangle = \langle T^*y, x\rangle$ untuk setiap $x$
--- lalu dengan mengonjugasikannya, $\langle Tx, y\rangle = \langle x,
T^*y\rangle$. Ketunggalannya membuat $T^*$ linear;
\[
\norm{T^*y} = \sup_{\norm x = 1}\abs{\langle T^*y, x\rangle}
= \sup_{\norm x=1}\abs{\langle y, Tx\rangle}
\leq \vertiii T\,\norm y,
\]
sehingga $\vertiii{T^*} \leq \vertiii T$, dan $T^{**} = T$ memberikan
kesamaannya. Untuk pergeserannya: $\langle Sx, y\rangle = \sum_{n\geq1}
x_n\bar y_{n+1} = \langle x, S^*y\rangle$ dengan $(S^*y)_n =
y_{n+1}$: yakni pergeseran mundurnya. Untuk kernel dan petanya: $T^*y = 0$ jika dan hanya jika
$\langle x, T^*y\rangle = 0$ untuk setiap $x$ jika dan hanya jika $\langle Tx,
y\rangle = 0$ untuk setiap $x$ jika dan hanya jika $y \perp \operatorname{im}T$:
sehingga $\ker T^* = (\operatorname{im}T)^\perp$; lalu dengan mengambil $\perp$ dan
memakai \cref{thm:b3:hilbert:decomposition},
$\overline{\operatorname{im}T} = (\ker T^*)^\perp$.
\end{solution}

\begin{solution}{exo:b3:hilbert:9}
Jika $a(u, \cdot) = \varphi$: maka untuk sembarang $w$,
\[
J(u + w) - J(u) = a(u, w) - \varphi(w) + \tfrac12a(w,w)
= \tfrac12a(w,w) \geq \tfrac\alpha2\norm w^2,
\]
yang positif sejati untuk $w \neq 0$: jadi $u$ satu-satunya pemberi minimumnya.
Sebaliknya, pada sebuah pemberi minimum fungsi $t \mapsto J(u + tw)$
(yakni polinomial kuadratik terhadap $t$) berturunan nol di
$0$: sehingga $a(u, w) - \varphi(w) = 0$ untuk setiap $w$. Untuk proyeksi
yang diturunkan kembali: bagi subruang tertutup $F$, terapkan Lax--Milgram pada
\omterm{def:b3:hilbert:inner}{ruang Hilbert} $F$ dengan $a(u,v) = \langle u, v\rangle$
(sehingga $M = \alpha = 1$) dan $\varphi(v) = \langle x, v\rangle$: diperoleh
tepat satu $p \in F$ dengan $\langle p, v\rangle = \langle x,
v\rangle$ untuk setiap $v \in F$, yakni $x - p \perp F$ --- dan menurut
kasus setangkupnya, $p$ meminimumkan $\frac12\norm v^2 - \langle
x, v\rangle = \frac12\norm{v - x}^2 - \frac12\norm x^2$ atas
$F$: yakni proyeksinya.
\end{solution}

\begin{solution}{exo:b3:hilbert:10}
Penormalannya: $\int h_n^2 = 2^j\cdot 2^{-j} = 1$.
Keortogonalannya: dua fungsi Haar yang berbeda entah
berpendukung (\omterm{def:b3:topology:interior}{interior}) saling lepas (sehingga hasil kalinya nol hampir di mana-mana), entah
pendukung yang lebih halus termuat di sebuah setengah selang tempat yang
lebih kasar bernilai konstan --- lalu integral hasil kalinya adalah
konstanta itu dikali $\int h_{\text{lebih halus}} = 0$; sedangkan terhadap $h_0
= \mathbf 1$, kembali $\int h_n = 0$. Untuk ketotalannya: rentang
$\{h_0, \dots, h_{2^J-1}\}$ terdiri atas fungsi tangga pada
kisi diadik berlangkah $2^{-J}$; kedua ruangnya berdimensi
$2^J$ dan fungsi Haarnya bebas (sebab ortonormal):
sehingga rentangnya adalah \emph{seluruh} fungsi tangga semacam itu. Lalu fungsi
tangga diadiknya padat di $L^2(\intcc01)$: sebab \omterm{def:b3:lebesgue:simple}{fungsi sederhananya}
padat (\cref{thm:b3:lp:density}(1)), himpunan \omterm{def:b3:lebesgue:measurable}{terukurnya}
dihampiri gabungan berhingga selang
(\cref{exo:b3:measure:7}), dan selangnya oleh yang diadik
(sebab ujungnya bergeser sejauh $\leq 2^{-J}$). Menurut
\cref{thm:b3:hilbert:parseval}, sistem Haarnya merupakan \omterm{def:b3:hilbert:onb}{basis
Hilbert}.
\end{solution}

\begin{solution}{exo:b3:hilbert:11}
(i) $\Rightarrow$ (ii): untuk proyeksi ortogonalnya,
$\langle Px, y\rangle = \langle Px, Py\rangle = \langle x,
Py\rangle$ (sisipkan penguraiannya $x = Px + (x - Px)$
dan seterusnya lalu bunuhlah suku silangnya). (ii) $\Rightarrow$ (iii):
$\norm{Px}^2 = \langle P^2x, x\rangle = \langle Px, x\rangle
\leq \norm{Px}\norm x$, sehingga $\vertiii P \leq 1$, dan $Pu = u$
pada petanya yang tak nol: jadi $= 1$. (iii) $\Rightarrow$ (i): $H =
\operatorname{im}P \oplus \ker P$ (secara aljabar, dari $P^2
= P$); lalu ambillah $u = Pu \in \operatorname{im}P$, $v \in \ker P$,
dan $t \in \R$: maka $\norm{P(u + tv)}^2 = \norm u^2$ haruslah $\leq
\norm{u + tv}^2 = \norm u^2 + 2t\operatorname{Re}\langle u,
v\rangle + t^2\norm v^2$ untuk setiap $t$, yang memaksa
$\operatorname{Re}\langle u, v\rangle = 0$ (bandingkan
suku linearnya saat $t \to 0^\pm$); lalu mengganti $v$ dengan $\iu v$
membunuh bagian imajinernya pula: sehingga $\operatorname{im}P \perp
\ker P$, yang tepat merupakan keortogonalan proyeksinya.
Contohnya: $P(x, y) = (x + y, 0)$ pada $\R^2$: maka $P^2 = P$,
dengan peta berupa sumbu-$x$, kernel berupa garis $y = -x$, dan
$\vertiii P = \sup\frac{\abs{x+y}}{\norm{(x,y)}} = \sqrt2$
(yang tercapai di $(1,1)/\sqrt2$): jadi proyeksi miring bernorma
$> 1$. (Untuk catatan, (ii) juga memberikan (i) secara langsung:
sebab $\ker P = \ker P^* = (\operatorname{im}P)^\perp$ menurut
\cref{exo:b3:hilbert:8}.)
\end{solution}

\begin{solution}{exo:b3:hilbert:12}
(a) Untuk $U$ yang uniter: $\norm{Ux - x}^2 = 2\norm x^2 -
2\operatorname{Re}\langle Ux, x\rangle$ dan $\norm{U^*x -
x}^2 = 2\norm x^2 - 2\operatorname{Re}\langle x, Ux\rangle$:
sehingga keduanya lenyap bersama, jadi $\ker(U - I) = \ker(U^* - I)$.
Lalu, dengan memakai $\ker T^* = (\operatorname{im}T)^\perp$
(\cref{exo:b3:hilbert:8}) untuk $T = U - I$ dan $T^* = U^* -
I$:
\[
\overline{\operatorname{im}(U - I)} = \bigl(\ker(U^* -
I)\bigr)^\perp = F^\perp .
\]

(b) Pada $F$: $U^kx = x$, sehingga $A_nx = x$. Untuk $x = (U - I)y$:
$A_nx = \frac1n(U^ny - y)$, yang bernorma $\leq \frac2n\norm y \to
0$. Sedangkan untuk $x$ pada \omterm{def:b3:topology:interior}{penutupnya}: diberikan $\varepsilon$, pilihlah $x' =
(U - I)y$ dengan $\norm{x - x'} < \varepsilon$; lalu karena
$\vertiii{A_n} \leq \frac1n\sum\vertiii{U^k} = 1$, maka
$\norm{A_nx} \leq \norm{A_n(x - x')} + \norm{A_nx'} \leq
\varepsilon + o(1)$.

(c) Uraikan $x = Px + (x - Px)$ dengan $Px \in F$ dan $x -
Px \in F^\perp = \overline{\operatorname{im}(U - I)}$ (bagian
(a)): sehingga $A_nx = Px + A_n(x - Px) \to Px + 0$.

(d) Pada basis Fouriernya $e_m(x) = \eu^{2\iu\pi mx}$: $Ue_m
= \eu^{2\iu\pi m\alpha}e_m$, sehingga $Ue_m = e_m$ jika dan hanya jika $m\alpha
\in \Z$ jika dan hanya jika $m = 0$ (sebab $\alpha$ irasional): jadi $F = \C\mathbf 1$
dan $Pf = \langle\mathbf 1, f\rangle\mathbf 1 = \int_0^1f$.
Teoremanya lalu berbunyi $\frac1n\sum_{k<n}f(\cdot + k\alpha) \to
\int_0^1f$ di $L^2(\R/\Z)$: sehingga rata-rata \omterm{def:b3:groups:action}{orbit} sebuah
rotasi irasional merata --- yakni bayangan $L^2$ dari
teorema pemerataan Weyl, yang diperoleh lewat geometri Hilbert belaka.
\end{solution}

\begin{solution}{pb:b3:hilbert:1}
\textbf{1.} Gram--Schmidt menjamin
$\operatorname{Vect}(p_0, \dots, p_n) =
\operatorname{Vect}(1, \dots, t^n) = \R_n[t]$ dan $p_n \perp
p_k$ (bila $k < n$), sehingga $p_n \perp \R_{n-1}[t]$. Sedangkan $p_k$, yang
berderajat naik sejati, bersifat bebas: jadi sebuah basis.

\textbf{2.} Polinomial $t\,p_n$ bersifat monik berderajat $n + 1$: uraikan
$t\,p_n = p_{n+1} + \sum_{k\leq n}c_kp_k$ dengan $c_k =
\langle p_k, tp_n\rangle/\norm{p_k}^2$. Untuk $k \leq n - 2$:
$\langle p_k, tp_n\rangle = \langle tp_k, p_n\rangle = 0$
(sebab berderajat $k + 1 < n$). Jadi $tp_n = p_{n+1} + a_np_n +
b_np_{n-1}$, yakni rekurensi yang dinyatakan itu, dengan
\[
b_n = \frac{\langle p_{n-1}, tp_n\rangle}{\norm{p_{n-1}}^2}
= \frac{\langle tp_{n-1}, p_n\rangle}{\norm{p_{n-1}}^2}
= \frac{\langle p_n + (\text{yang lebih rendah}),\
p_n\rangle}{\norm{p_{n-1}}^2}
= \frac{\norm{p_n}^2}{\norm{p_{n-1}}^2} > 0 .
\]

\textbf{3.} Misalkan $t_1 < \dots < t_m$ titik \omterm{def:b3:topology:interior}{interior}
$I$ yang $p_n$-nya berganti tanda, dan $q = \prod_{i\leq m}(t -
t_i)$ (dengan $q = 1$ bila $m = 0$). Maka $p_nq$ bertanda tetap
pada $I$ dan tak nol hampir di mana-mana: sehingga $\int_Ip_nq\,w \neq 0$. Jika $m <
n$, ini bertentangan dengan $p_n \perp \R_{n-1}[t]$. Jadi $m = n$:
sehingga $p_n$ mempunyai $n$ akar \omterm{def:b3:topology:interior}{interior} yang berbeda (sebab ia berakar paling banyak $n$
seluruhnya).

\textbf{4.} Polinomial $(t^2 - 1)^n$ berderajat $2n$; lalu $n$ turunan
menyisakan derajat $n$, dengan koefisien utama
$\frac{(2n)(2n-1)\cdots(n+1)}{2^nn!} =
\frac{(2n)!}{2^n(n!)^2}$. Untuk $\deg Q < n$, integralkan secara parsial
$n$ kali: maka semua suku batasnya memuat sebuah turunan berorde
$< n$ dari $(t^2-1)^n$, yang lenyap di $\pm1$ (sebab akar berorde
$n$); dan setelah $n$ langkah integrannya membawa $Q^{(n)} = 0$.

\textbf{5.} Dengan $u = (t^2 - 1)^n$:
\[
(2^nn!)^2\norm{P_n}^2 = \int_{-1}^1(u^{(n)})^2
= (-1)^n\int_{-1}^1 u\,u^{(2n)}
= (2n)!\int_{-1}^1(1 - t^2)^n\dd t ,
\]
(sebab $u^{(2n)} = (2n)!$; sedangkan suku batasnya lenyap seperti pada pertanyaan 4).
Dan $\int_{-1}^1(1-t^2)^n\dd t = B(\tfrac12, n+1) =
\frac{\Gamma(\frac12)\Gamma(n+1)}{\Gamma(n + \frac32)} =
\frac{2\cdot4^n(n!)^2}{(2n+1)!}$
(\cref{exo:b3:product:8}). Digabungkan: $\norm{P_n}^2 =
\frac{2}{2n + 1}$.

\textbf{6.} Polinomialnya padat terhadap $\norm\cdot_\infty$ di
$\mathcal C(\intcc{-1}1)$ (menurut Weierstrass,
\cref{cor:b3:complete:weierstrass}), fungsi \omterm{def:b3:topology:continuity}{kontinunya}
padat-$L^2$ (\cref{thm:b3:lp:density}), dan
$\norm\cdot_2 \leq \sqrt2\norm\cdot_\infty$: sehingga rentang polinomialnya
total, jadi $P_n$ yang ternormalkan membentuk \omterm{def:b3:hilbert:onb}{basis Hilbert}.
Untuk uraian $\abs t$: koefisiennya terhadap $P_0$ adalah
$\frac{\langle P_0, \abs t\rangle}{\norm{P_0}^2} = \frac12$;
terhadap $P_1$: $0$ (menurut paritasnya); dan terhadap $P_2$:
$\frac{\int_{-1}^1\abs t\,\frac{3t^2-1}2\dd t}{2/5} =
\frac{1/4}{2/5} = \frac58$. Jadi penghampiran kuadratik terbaiknya:
\[
\abs t \approx \frac12 + \frac58\,P_2(t) = \frac{3}{16} +
\frac{15}{16}\,t^2 .
\]

\textbf{7.} Dari $\frac{\dd^{n+1}}{\dd t^{n+1}}\eu^{-t^2} =
\frac{\dd^n}{\dd t^n}(-2t\,\eu^{-t^2})$ dan Leibniz: $H_{n+1}
= 2tH_n - H_n'$; lalu induksinya memberikan derajat $n$ dan koefisien
utama $2^n$. Untuk $m < n$, integralkan secara parsial $n$ kali
pada $\int H_m H_n\eu^{-t^2} = (-1)^n\int H_m\,\bigl(\eu^{-t^2}
\bigr)^{(n)}$: maka suku batasnya (yakni polinomial $\times$
$\eu^{-t^2}$) lenyap di $\pm\infty$, menyisakan $\int
H_m^{(n)}\,\eu^{-t^2} = 0$. Sedangkan untuk $m = n$: $H_n^{(n)} = 2^nn!$,
sehingga $\norm{H_n}_w^2 = 2^nn!\int\eu^{-t^2} = 2^nn!\sqrt\pi$.

\textbf{8.} Misalkan $f \in L^2(\R, \eu^{-t^2}\dd t)$ ortogonal
terhadap setiap polinomial, dan $g = f\eu^{-t^2}$. Maka $g \in
L^1$: sebab $\int\abs f\eu^{-t^2} \leq \bigl(\int\abs
f^2\eu^{-t^2}\bigr)^{1/2}\bigl(\int\eu^{-t^2}\bigr)^{1/2}$
(lewat Cauchy--Schwarz). Untuk $\xi \in \R$, uraikan $\eu^{-\iu\xi t}$:
maka jumlah parsialnya terdominasi sebab
\[
\sum_k\frac{\abs\xi^k}{k!}\int\abs f\,\abs t^k\eu^{-t^2}\dd t
\leq \Bigl(\int \abs f^2\eu^{-t^2}\Bigr)^{1/2}
\sum_k\frac{\abs\xi^k}{k!}\Bigl(\int
t^{2k}\eu^{-t^2}\Bigr)^{1/2} < \infty
\]
(sebab deret terakhirnya konvergen: $\int t^{2k}\eu^{-t^2} =
\Gamma(k+\frac12) \leq k!\,\sqrt\pi$, sehingga sukunya bernilai
$O(\abs\xi^k/\sqrt{k!})$). Lalu pengintegralan suku demi suku
(\cref{cor:b3:lebesgue:additivity} yang diterapkan pada deret
mutlaknya, lalu Fubini untuk deret) memberikan
\[
\int_\R g(t)\,\eu^{-\iu\xi t}\dd t
= \sum_k\frac{(-\iu\xi)^k}{k!}\int f(t)\,t^k\,\eu^{-t^2}\dd t
= 0 ,
\]
sebab setiap integralnya bertipe $\langle t^k, f\rangle_w$ yang $= 0$. Menurut
keinjektifan transformasi Fourier yang diterima
(\cref{ch:b3:fouriertransform}), $g = 0$ hampir di mana-mana, sehingga $f = 0$
hampir di mana-mana: jadi keluarga Hermitenya (yang rentangnya adalah polinomialnya) bersifat
total.

\textbf{9.} Untuk kepersisannya sampai derajat $n - 1$: bagi $P$ semacam itu, $P =
\sum_iP(t_i)\ell_i$ secara persis, sehingga $\int Pw = \sum_iP(t_i)\int
\ell_iw = Q(P)$. Untuk derajat $\leq 2n - 1$: bagilah $P = qp_n + r$
dengan $\deg q \leq n - 1$ dan $\deg r \leq n-1$; maka $\int Pw = \int
qp_nw + \int rw = 0 + Q(r)$ (sebab $p_n \perp \R_{n-1}[t]$), sedangkan
$Q(P) = \sum_iw_i\bigl(q(t_i)\,p_n(t_i) + r(t_i)\bigr) = Q(r)$
sebab simpulnya adalah akar $p_n$. Jadi sama.

\textbf{10.} Polinomial $\ell_i^2$ berderajat $2n - 2 \leq 2n - 1$ dan
$\ell_i^2(t_j) = \delta_{ij}$: sehingga $0 < \int\ell_i^2w = Q(\ell_i^2)
= w_i$. Lalu Polya (\cref{exo:b3:banach:9}, yang dipindahkan ke $I$ dengan
bobotnya): syarat (i) berlaku --- sebab setiap polinomialnya diintegralkan
persis sekali karena $2n - 1 \geq$ derajatnya; sedangkan syarat (ii):
$\sum_i\abs{w_{i}} = \sum_iw_i = Q(\mathbf 1) = \int_Iw$,
yang terbatas: sehingga $Q_n(f) \to \int fw$ untuk setiap $f \in \mathcal C(I)$
dengan $I$ \omterm{def:b3:topology:compact}{kompak}.

\textbf{11.} Yang monik $p_2 = t^2 - \frac13$ (dari
\cref{exo:b3:hilbert:4}): sehingga simpulnya $\pm\frac1{\sqrt3}$. Bobotnya:
$\ell_1(t) = \frac{t - \frac1{\sqrt3}}{-\frac2{\sqrt3}}$, dan
$w_1 = \int_{-1}^1\ell_1 = 1$; lalu menurut kesetangkupannya $w_2 = 1$.
Kepersisannya: $\int 1 = 2 = 1 + 1$; $\int t = 0 =
-\frac1{\sqrt3} + \frac1{\sqrt3}$; $\int t^2 = \frac23 =
\frac13 + \frac13$; $\int t^3 = 0$. Sedangkan aturan trapesium dua
titiknya (dengan simpul $\pm1$ dan bobot $1, 1$) persis hanya sampai derajat
$1$: sebab pada $t^2$ ia mengembalikan $2$ alih-alih $\frac23$. Jadi dengan biaya yang sama,
dua derajat kepersisan tambahan: itulah hasil simpul yang
ortogonal.

\textbf{12.} Dari $\cos(n{+}1)\theta + \cos(n{-}1)\theta =
2\cos\theta\cos n\theta$: $T_{n+1} = 2tT_n - T_{n-1}$ dengan
$T_0 = 1$ dan $T_1 = t$; lalu induksinya memberikan polinomial berderajat
$n$ dengan koefisien utama $2^{n-1}$ (bila $n \geq 1$).
Dengan menyubstitusikan $t = \cos\theta$ (sehingga $w(t)\dd t \mapsto
\dd\theta$): $\langle T_m, T_n\rangle_w =
\int_0^\pi\cos m\theta\cos n\theta\,\dd\theta = 0$ untuk $m
\neq n$, $= \pi$ untuk $m = n = 0$, dan $= \frac\pi2$ selainnya
(lewat hasil kali ke jumlah). Derajat dan keortogonalan berpasangannya
mengenali $T_n$ dengan keluaran Gram--Schmidt sampai
skalarnya; sedangkan uraian Chebyshev atas $f$ tepat merupakan
deret kosinus Fourier dari $\theta \mapsto f(\cos\theta)$.

\textbf{13.} Berlaku $T_n(t) = 0$ jika dan hanya jika $\cos n\theta = 0$ jika dan hanya jika
$\theta = \frac{(2k-1)\pi}{2n}$: sehingga $n$ akarnya berbeda, yakni
$t_k = \cos\frac{(2k-1)\pi}{2n} \in \intoo{-1}1$. Ekstremumnya:
$\abs{T_n} \leq 1$ pada $\intcc{-1}1$, dengan $T_n(s_j) =
(-1)^j$ di $n + 1$ titik $s_j = \cos\frac{j\pi}n$:
yakni osilasi setara yang \omterm{prop:b3:galois:perfect}{sempurna}.

\textbf{14.} Polinomial $2^{1-n}T_n$ bersifat monik dengan norma-sup $2^{1-n}$.
Jika sebuah $P$ monik berderajat $n$ mempunyai $\sup\abs P < 2^{1-n}$, maka
selisih $D = 2^{1-n}T_n - P$ akan berderajat $\leq n -
1$ (sebab suku utamanya saling meniadakan) padahal berganti tanda di $s_0 >
\dots > s_n$ (sebab di sana $2^{1-n}T_n = \pm2^{1-n}$ mendominasi
$P$): jadi setidaknya $n$ akar --- sehingga $D \equiv 0$, yang bertentangan.
Untuk ketunggalannya pada kesamaannya, $D$ yang sama memenuhi
$(-1)^jD(s_j) \geq 0$; dan polinomial tak nol berderajat $\leq
n-1$ tak dapat mempunyai $n$ kendala ekstremal yang berganti tanda
secara lemah tanpa $n$ akar yang tercacah dengan benar (sebab jika
$D(s_j) = 0$ untuk suatu $s_j$ \omterm{def:b3:topology:interior}{interior}, maka akar itu rangkap
dalam pencacahannya karena $D$ mempertahankan tanda secara lokal): sehingga lagi-lagi $D
\equiv 0$.

\textbf{15.} Rumus galat Lagrange (lewat Rolle, Tahun~2)
memberikan $f - L_nf = \frac{f^{(n)}(\xi_t)}{n!}\,\omega(t)$, sehingga
galat seragamnya paling banyak $\frac{\norm{f^{(n)}}_\infty}
{n!}\,\sup\abs\omega$, sedangkan $\omega$ monik berderajat $n$:
sehingga menurut pertanyaan 14, $\sup_{\intcc{-1}1}\abs\omega \geq 2^{1-n}$
dengan kesamaannya jika dan hanya jika simpulnya akar Chebyshev. Karena itu
batas optimalnya $\norm{f - L_nf}_\infty \leq
\frac{\norm{f^{(n)}}_\infty}{2^{n-1}n!}$. Dengan simpul berjarak
sama, $\sup\abs\omega$ jauh lebih besar secara eksponensial di dekat
ujungnya, dan menginterpolasi bahkan $\frac1{1 + 25t^2}$
divergen di sana saat $n \to \infty$ (yakni gejala Runge);
sedangkan simpul Chebyshev adalah obatnya.

\textbf{16.} Dengan menurunkan $T_n(\cos\theta) = \cos
n\theta$: $T_n'(\cos\theta) = \frac{n\sin n\theta}
{\sin\theta}$, yang menuju $n^2$ saat $\theta \to 0$ dan
$(-1)^{n+1}n^2$ saat $\theta \to \pi$: sehingga $\abs{T_n'(\pm1)} =
n^2$. Sedangkan di titik \omterm{def:b3:topology:interior}{interiornya}, $\abs{T_n'(t)} \leq
\frac{n}{\sqrt{1 - t^2}} = O(n)$: jadi peledakan kuadratiknya
hidup hanya di tepinya (yakni batas \omterm{def:b3:topology:interior}{interior} Bernstein lawan
batas global Markov).

\textbf{17.} Misalkan $\theta_k = \frac{(2k-1)\pi}{2n}$ dan $S_j
= \sum_{k=1}^n\cos(j\theta_k)$ untuk $1 \leq j \leq n-1$. Maka
\[
S_j = \operatorname{Re}\Bigl[\eu^{\iu j\pi/2n}
\sum_{k=0}^{n-1}\eu^{\iu jk\pi/n}\Bigr]
= \operatorname{Re}\Bigl[\eu^{\iu j\pi/2n}\,
\frac{\eu^{\iu j\pi} - 1}{\eu^{\iu j\pi/n} - 1}\Bigr] .
\]
Untuk $j$ yang genap pembilangnya lenyap: sehingga $S_j = 0$. Sedangkan untuk $j$ yang gasal
pembilangnya $-2$, dan $\eu^{\iu j\pi/n} - 1 =
\eu^{\iu j\pi/2n}\cdot2\iu\sin\frac{j\pi}{2n}$, sehingga seluruh
ungkapannya menjadi $\frac{-2}{2\iu\sin(j\pi/2n)} =
\frac{\iu}{\sin(j\pi/2n)}$: yang murni imajiner, jadi $S_j = 0$
lagi. Karena itu aturan berbobot sama $\frac\pi n\sum_kf(t_k)$
mengintegralkan $T_0$ (sebab $\sum w_i = \pi = \int w$) dan membunuh $T_1,
\dots, T_{n-1}$ persis seperti $\int T_jw = 0$: sehingga ia persis
sampai derajat $n - 1$. Sedangkan bobot yang persis sampai derajat $n-1$ pada simpul
tertentu bersifat tunggal (lewat basis Lagrange): jadi semua bobot Gaussnya
$\frac\pi n$. Untuk $n = 3$: simpulnya $\pm\frac{\sqrt3}2, 0$ dan
\[
\int_{-1}^1\frac{f(t)}{\sqrt{1 - t^2}}\,\dd t \approx
\frac\pi3\Bigl[f\Bigl(\tfrac{\sqrt3}2\Bigr) + f(0) +
f\Bigl(-\tfrac{\sqrt3}2\Bigr)\Bigr],
\]
persis sampai derajat $5$.

\textbf{18.} Untuk $P$ monik berderajat $n$: $P = p_n + r$
dengan $r \in \R_{n-1}[t]$, dan $p_n \perp \R_{n-1}[t]$
(pertanyaan 1), sehingga $\norm P^2 = \norm{p_n}^2 + \norm r^2 \geq
\norm{p_n}^2$, dengan kesamaannya jika dan hanya jika $r = 0$: jadi $p_n$ merupakan
sisa proyeksi ortogonal $t^n$ ke
$\R_{n-1}[t]^\perp$, yakni polinomial monik terdekat dengan
subruang yang harus dihindarinya. Sedangkan $2^{1-n}T_n$ milik Chebyshev menjawab
pertanyaan yang sama untuk norma supnya: yakni simpangan terkecil dari
nol, sekali di $L^2(w)$, sekali di $L^\infty$.

\textbf{19.} Tulis $K_n(x, y) =
\sum_{k=0}^n\frac{p_k(x)p_k(y)}{h_k}$. Basisnya $n = 0$: $(x -
y)\frac1{h_0} = \frac{p_1(x)\cdot1 - 1\cdot p_1(y)}{h_0}$
sebab $p_1 = t - a_0$. Langkahnya: dengan mengandaikan identitasnya untuk $n -
1$,
\[
(x - y)\,K_n(x,y) = \frac{p_n(x)p_{n-1}(y) -
p_{n-1}(x)p_n(y)}{h_{n-1}} +
\frac{(x - y)\,p_n(x)p_n(y)}{h_n} ;
\]
substitusikan $x\,p_n(x) = p_{n+1}(x) + a_np_n(x) +
b_np_{n-1}(x)$ dan $y\,p_n(y) = p_{n+1}(y) + a_np_n(y) +
b_np_{n-1}(y)$ pada suku keduanya: maka sumbangan $a_n$-nya
saling meniadakan, dan sumbangan $b_n = \frac{h_n}{h_{n-1}}$-nya
meniadakan suku induksinya; sehingga yang bertahan adalah
$\frac{p_{n+1}(x)p_n(y) - p_n(x)p_{n+1}(y)}{h_n}$. Sedangkan
bentuk konfluennya menyusul dengan membiarkan $y \to x$ (sebab kedua ruasnya
polinomial terhadap $y$).

\textbf{20.} Bentuk konfluennya memberikan $p_{n+1}'p_n -
p_n'p_{n+1} = h_n\sum_{k\leq n}\frac{p_k^2}{h_k} \geq
\frac{h_n}{h_0} > 0$ di mana-mana. Pada sebuah akar $x_0$ dari
$p_{n+1}$: $p_{n+1}'(x_0)\,p_n(x_0) > 0$, sehingga $p_n(x_0) \neq
0$ (jadi tak ada akar bersama). Di antara akar berurutan $x_0 < x_1$
dari $p_{n+1}$ (yang semuanya sederhana, Bagian I), $p_{n+1}'$ bertanda
berlawanan, sehingga demikian pula $p_n$: jadi sebuah akar $p_n$ terletak di setiap
$n$ celahnya --- dan itu menghabiskan $n$ akarnya: yakni
keberselang-selingannya.

\textbf{21.} Dengan menguraikan $D_n(t) = \det(tI_n - J_n)$ sepanjang
baris terakhirnya: $D_n = (t - a_{n-1})D_{n-1} - b_{n-1}D_{n-2}$,
dengan $D_0 = 1$ dan $D_1 = t - a_0$: yakni rekurensi dan benih
$p_n$ yang monik, sehingga $D_n = p_n$. Jadi akar $p_n$ adalah
nilai eigen $J_n$ yang setangkup: yang real, dan sederhana menurut
pertanyaan 19 --- sehingga kuadratur Gauss adalah teori spektral sebuah
matriks tridiagonal yang menyamar, yakni bayangan berdimensi berhingga
\cref{ch:b3:spectral}.

\textbf{22.} Kamusnya:
\begin{center}
\begin{tabular}{l|c|c|c}
 & Legendre & Hermite & Chebyshev\\
\hline
selang & $\intcc{-1}1$ & $\R$ & $\intcc{-1}1$\\
bobot & $1$ & $\eu^{-t^2}$ & $(1-t^2)^{-1/2}$\\
rumus & Rodrigues & $(-1)^n\eu^{t^2}
\frac{\dd^n}{\dd t^n}\eu^{-t^2}$ & $\cos(n\arccos t)$\\
norma$^2$ & $\frac2{2n+1}$ & $2^nn!\sqrt\pi$ & $\pi,
\frac\pi2$\\
habitat & kuadratur & kalkulus Gauss & minimaks\\
\end{tabular}
\end{center}
(masing-masing dengan rekurensi tiga sukunya: yakni bentuk umumnya untuk
Legendre, $H_{n+1} = 2tH_n - 2nH_{n-1}$, dan $T_{n+1} = 2tT_n -
T_{n-1}$). Teori umumnya menyediakan apa yang tak ditunjukkan satu keluarga
mana pun: yakni kerealan dan keberselang-selingan akarnya, kepositifan
bobot kuadraturnya, serta keberadaan rekurensinya dan
Christoffel--Darboux belaka --- semuanya akibat keortogonalannya
saja, yang seragam terhadap bobotnya.

\textbf{23.} Keberadaannya: pemetaan linear $\R_{2n-1}[t] \to
\R^{2n}$, $P \mapsto (P(t_1), P'(t_1), \dots, P(t_n),
P'(t_n))$, bersifat injektif (sebab $P$ di kernelnya berakar rangkap $n$
kali dan berderajat $\leq 2n - 1$, jadi $P = 0$) antara ruang
berdimensi sama $2n$: sehingga bijektif. Untuk galat titik demi titiknya: tetapkan $t$
yang bukan simpul lalu pilihlah $K$ sedemikian sehingga $g(s) = f(s) - Hf(s) -
K\,p_n(s)^2$ lenyap di $s = t$. Maka $g$ lenyap di
$n + 1$ titik berbeda $t, t_1, \dots, t_n$, dan $g'$
lenyap di setiap $t_i$ pula (sebab baik $f - Hf$ maupun $p_n^2$ berakar
rangkap di sana). Lalu Rolle memberikan $n$ akar $g'$ yang terletak sejati
di antara akar berurutan $g$ --- yang berbeda dari simpulnya
--- sehingga $g'$ mempunyai $2n$ akar berbeda; lalu menerapkan Rolle $2n - 1$
kali lagi menghasilkan $\xi_t$ dengan $g^{(2n)}(\xi_t) = 0$.
Karena $\deg Hf \leq 2n - 1$ dan $p_n^2$ monik berderajat
$2n$, maka $g^{(2n)} = f^{(2n)} - K\,(2n)!$, sehingga $K =
f^{(2n)}(\xi_t)/(2n)!$ --- dan identitasnya trivial di
simpulnya. Untuk pengintegralannya: $Q_n(f) = Q_n(Hf)$ (sebab $Hf$ mencocokkan $f$
di simpulnya) dan $Q_n(Hf) = \int Hf\,w$ menurut kepersisannya sampai
derajat $2n - 1$ (pertanyaan 9), sehingga galat kuadraturnya adalah
$\int(f - Hf)\,w$. Lalu dengan $m, M$ sebagai ekstremum $f^{(2n)}$ pada
$I$, identitas titik demi titiknya menjepit
\[
\frac{m\,h_n}{(2n)!} \;\leq\; \int_I(f - Hf)\,w
\;\leq\; \frac{M\,h_n}{(2n)!} ,
\]
lalu teorema nilai antara yang diterapkan pada $f^{(2n)}$ yang
\omterm{def:b3:topology:continuity}{kontinu} melahirkan $\xi$. (Untuk Legendre dengan $n = 2$: $h_2
= \int_{-1}^1(t^2 - \frac13)^2\dd t = \frac8{45}$, sehingga
galatnya $f^{(4)}(\xi)/135$.)

\textbf{24.} Kernelnya mereproduksi $\R_{n-1}[t]$: sebab menguraikan
$q = \sum_k\frac{\langle p_k, q\rangle}{h_k}p_k$ memberikan
$\int_I K_n(t_i, t)\,q(t)\,w(t)\dd t = q(t_i)$ untuk setiap
$q$ berderajat $\leq n - 1$. Ambillah $q = \ell_i$: maka ruas kirinya
sama dengan $\ell_i(t_i) = 1$. Padahal $t \mapsto K_n(t_i,
t)\,\ell_i(t)$ merupakan polinomial berderajat $\leq (n - 1) + (n
- 1) = 2n - 2$, yang $Q_n$-nya persis (pertanyaan 9), dan ia
lenyap di setiap simpul $t_j \neq t_i$ (lewat faktor $\ell_i$), sehingga
\[
1 = \int_I K_n(t_i, t)\,\ell_i(t)\,w(t)\dd t
= w_i\,K_n(t_i, t_i)
= w_i\sum_{k=0}^{n-1}\frac{p_k(t_i)^2}{h_k} .
\]
Jumlahnya $> 0$ (sebab suku $k = 0$-nya adalah $1/h_0 > 0$): sehingga diperoleh
rumus yang dinyatakan itu, beserta kepositifannya lagi. Pemeriksaannya ($n = 2$,
Legendre): $p_0 = 1$, $h_0 = 2$, $p_1 = t$, $h_1 = \frac23$;
dan di $t_i = \pm\frac1{\sqrt3}$,
\[
K_2(t_i, t_i) = \frac12 + \frac{1/3}{2/3} = 1,
\qquad w_i = 1,
\]
seperti yang ditemukan pada pertanyaan 11.

\textbf{25.} Dengan menyubstitusikan $t = \cos\theta$, integralnya menjadi
$\int_0^\pi\cos^6\theta\,\dd\theta =
\pi\,\frac{5\cdot3\cdot1}{6\cdot4\cdot2} = \frac{5\pi}{16}$
(lewat Wallis, \cref{exo:b3:product:8}). Sedangkan aturan
Chebyshev--Gauss dengan $n = 3$ (pertanyaan 17) bersimpul
$\cos\frac\pi6 = \frac{\sqrt3}2$, $\cos\frac\pi2 = 0$, dan
$\cos\frac{5\pi}6 = -\frac{\sqrt3}2$ serta berbobot sama
$\frac\pi3$:
\[
Q_3(t^6) = \frac\pi3\Bigl(2\cdot\Bigl(\frac{\sqrt3}2
\Bigr)^{6}\Bigr) = \frac\pi3\cdot\frac{54}{64}
= \frac{9\pi}{32},
\qquad
\frac{5\pi}{16} - \frac{9\pi}{32} = \frac\pi{32} .
\]
Ramalannya: \omterm{pb:b3:hilbert:1}{polinomial ortogonal} monik berderajat $3$ adalah
$2^{-2}T_3 = t^3 - \frac34t$, dengan $h_3 =
\frac1{16}\norm{T_3}_w^2 = \frac1{16}\cdot\frac\pi2 =
\frac\pi{32}$; dan $f = t^6$ mempunyai $f^{(6)} = 720 =
6!$ yang konstan, sehingga pertanyaan 23 memberikan galat $\frac{6!}{6!}\,h_3 =
\frac\pi{32}$ --- dan karena tak ada kebergantungan pada $\xi$ yang tersisa,
rumusnya terpaksa persis, dan memang demikian.

\end{solution}
